\section{HHLアルゴリズム}
\label{sec:hhl-algorithm}

HHLアルゴリズムは，量子位相推定，固有値に応じた制御回転，逆量子位相推定，
補助量子ビットの測定から構成される\cite{hhl2009}．
以下では，$A$の固有値が既知の範囲に収まるようにスケーリングされ，
$e^{\mathrm{i}At}$を量子回路として実行できるものとする．

\subsection{レジスタの準備}
\label{subsec:hhl-register-preparation}

HHLアルゴリズムでは，次の三つのレジスタを使用する．

\begin{enumerate}[(1)]
  \item 線形方程式の右辺$\ket{b}$を保持する状態レジスタ
  \item 固有値の近似値を保持する位相レジスタ
  \item 固有値の逆数を振幅へ反映するための補助量子ビット
\end{enumerate}

初期状態を
\begin{align*}
  \ket{b}\ket{0}^{\otimes r}\ket{0}_{a}
\end{align*}
とする．
ここで，$r$は位相レジスタの量子ビット数，下付きの$a$は補助量子ビットを表す．
式\eqref{eq:b-eigenbasis-expansion}を用いると，状態レジスタは
\begin{align*}
  \ket{b}
  =
  \sum_j\beta_j\ket{u_j}
\end{align*}
と表される．

\subsection{量子位相推定}
\label{subsec:hhl-quantum-phase-estimation}

行列$A$から定まるユニタリ演算子を
\begin{equation}
  U_A(t)
  \coloneq
  e^{\mathrm{i}At}
  \label{eq:hamiltonian-simulation-unitary}
\end{equation}
と定義する．
$A\ket{u_j}=\lambda_j\ket{u_j}$であるため，
\begin{align*}
  U_A(t)\ket{u_j}
  =
  e^{\mathrm{i}\lambda_jt}\ket{u_j}
\end{align*}
となる．
したがって，$\ket{u_j}$を入力とする量子位相推定により，
位相$\lambda_jt$から固有値$\lambda_j$を近似できる．

量子位相推定を状態$\ket{b}$へ線形に作用させると，理想的には
\begin{equation}
  \sum_j\beta_j\ket{u_j}\ket{0}^{\otimes r}
  \longmapsto
  \sum_j\beta_j\ket{u_j}\ket{\widetilde{\lambda}_j}
  \label{eq:hhl-phase-estimation}
\end{equation}
となる．
ここで，$\widetilde{\lambda}_j$は有限個の量子ビットで表された$\lambda_j$の近似値である．
位相レジスタを増やすと固有値の分解能は高くなるが，必要な制御演算と回路深さも増加する．

\subsection{固有値の逆数を反映する制御回転}
\label{subsec:hhl-controlled-rotation}

次に，位相レジスタの値$\widetilde{\lambda}_j$に応じて補助量子ビットを回転させる．
定数$C$を
\begin{align*}
  0<C\leq\min_j|\lambda_j|
\end{align*}
となるように選び，次の変換を行う．
\begin{equation}
  \ket{\widetilde{\lambda}_j}\ket{0}_{a}
  \longmapsto
  \ket{\widetilde{\lambda}_j}
  \left(
    \sqrt{1-\frac{C^2}{\widetilde{\lambda}_j^2}}\ket{0}_{a}
    +\frac{C}{\widetilde{\lambda}_j}\ket{1}_{a}
  \right)
  \label{eq:hhl-controlled-rotation}
\end{equation}
これにより，補助量子ビットが$\ket{1}_{a}$である成分の振幅へ
$1/\widetilde{\lambda}_j$が反映される．

量子位相推定後の状態へ式\eqref{eq:hhl-controlled-rotation}を作用させると
\begin{align*}
  \sum_j\beta_j\ket{u_j}\ket{\widetilde{\lambda}_j}
  \left(
    \sqrt{1-\frac{C^2}{\widetilde{\lambda}_j^2}}\ket{0}_{a}
    +\frac{C}{\widetilde{\lambda}_j}\ket{1}_{a}
  \right)
\end{align*}
を得る．
小さな固有値ほど回転角への影響が大きいため，固有値推定の誤差と条件数がこの処理の精度を左右する．

\subsection{逆量子位相推定と補助量子ビットの測定}
\label{subsec:hhl-uncomputation-and-measurement}

制御回転の後，量子位相推定の逆回路を適用し，位相レジスタを$\ket{0}^{\otimes r}$へ戻す．
固有値の近似誤差を無視すると，状態レジスタと補助量子ビットは
\begin{align*}
  \sum_j\beta_j\ket{u_j}
  \left(
    \sqrt{1-\frac{C^2}{\lambda_j^2}}\ket{0}_{a}
    +\frac{C}{\lambda_j}\ket{1}_{a}
  \right)
\end{align*}
となる．

補助量子ビットを計算基底で測定し，結果$1$を得た場合，状態レジスタは
\begin{align*}
  C\sum_j\frac{\beta_j}{\lambda_j}\ket{u_j}
  =
  CA^{-1}\ket{b}
\end{align*}
に比例する状態へ射影される．
これを規格化すると，式\eqref{eq:qlsp-solution-state}の$\ket{x}$が得られる．
理想的な場合に補助量子ビットから$1$を得る確率は
\begin{equation}
  p_1
  \coloneq
  C^2\sum_j\frac{|\beta_j|^2}{\lambda_j^2}
  =
  C^2\|A^{-1}\ket{b}\|_2^2
  \label{eq:hhl-success-probability}
\end{equation}
である．
測定結果が$0$であれば解状態は得られないため，状態準備から処理をやり直す必要がある．

\subsection{処理手順のまとめ}
\label{subsec:hhl-procedure-summary}

HHLアルゴリズムの処理は，次の手順にまとめられる．

\begin{enumerate}[(1)]
  \item 右辺ベクトルを規格化し，状態レジスタへ$\ket{b}$として準備する．
  \item $e^{\mathrm{i}At}$を用いた量子位相推定により，$A$の固有値を位相レジスタへ書き込む．
  \item 固有値の逆数に比例する振幅を，制御回転によって補助量子ビットへ書き込む．
  \item 逆量子位相推定によって位相レジスタの計算を巻き戻す．
  \item 補助量子ビットを測定し，結果$1$に条件づけて解状態を得る．
\end{enumerate}

このように，HHLアルゴリズムは逆行列$A^{-1}$を古典的に作成するのではなく，
固有値ごとに振幅を変換することで$A^{-1}\ket{b}$に比例する量子状態を生成する．
